% =====================================================================
%  دليل الاستخدام: الألوان والرموز والصناديق
% =====================================================================
\chapter{دليل الاستخدام}

\section{طريقة الشرح}
نمشي على الشاطبية \textbf{بيتًا بيتًا}. لكل بيت:
\begin{enumerate}
\item \textbf{نص البيت} في إطار ذهبي، ورقمه بجانبه (بحسب ترقيم المتن المضبوط).
\item \textbf{معاني الكلمات} الصعبة فقط.
\item \textbf{المعنى}: شرح مختصر سهل.
\item في أبواب القراءات: \textbf{صندوق القراءات} يبين من قرأ بكذا، ثم \textbf{القاعدة} إن كانت تحتاج تفسيرًا، ثم \textbf{أمثلة من القرآن}، ثم \textbf{التحريرات} إن وُجدت.
\end{enumerate}
وفي آخر كل باب \textbf{خلاصة} تجمع أحكامه في جدول.

\section{ألوان القراء}
كل قارئ له لون ثابت في المستند كله، وراوياه يأخذان لونه نفسه. فإذا رأيت اللون الأحمر فهو \Asim{} أو أحد راوييه، وهكذا.

\begin{center}
\renewcommand{\arraystretch}{1.5}
\begin{tabular}{c c c c c}
\toprule
\textbf{القارئ} & \textbf{رمزه} & \textbf{الراوي الأول} & \textbf{الراوي الثاني} & \textbf{بلده}\\
\midrule
\Nafi & \rzl{Nafi}{أ} & \Qalun{} \rzl{Nafi}{ب} & \Warsh{} \rzl{Nafi}{ج} & المدينة\\
\IbnKathir & \rzl{IbnKathir}{د} & \Bazzi{} \rzl{IbnKathir}{هـ} & \Qunbul{} \rzl{IbnKathir}{ز} & مكة\\
\AbuAmr & \rzl{AbuAmr}{ح} & \qname{AbuAmr}{الدوري} \rzl{AbuAmr}{ط} & \Susi{} \rzl{AbuAmr}{ي} & البصرة\\
\IbnAmir & \rzl{IbnAmir}{ك} & \Hisham{} \rzl{IbnAmir}{ل} & \IbnDhakwan{} \rzl{IbnAmir}{م} & الشام\\
\Asim & \rzl{Asim}{ن} & \Shuba{} \rzl{Asim}{ص} & \Hafs{} \rzl{Asim}{ع} & الكوفة\\
\Hamza & \rzl{Hamza}{ف} & \Khalaf{} \rzl{Hamza}{ض} & \Khallad{} \rzl{Hamza}{ق} & الكوفة\\
\Kisai & \rzl{Kisai}{ر} & \AbulHarith{} \rzl{Kisai}{س} & \qname{Kisai}{الدوري} \rzl{Kisai}{ت} & الكوفة\\
\bottomrule
\end{tabular}
\end{center}

\begin{tanbih}[طريقة حفظ الرموز]
الحروف مرتبة على «أَبْجَدْ هَوَّزْ حُطِّي كَلَمُنْ سَعْفَصْ قَرَشَتْ»، كل ثلاثة أحرف لقارئ وراوييه بالترتيب:
\textbf{\rzl{Nafi}{أبج}} \textbf{\rzl{IbnKathir}{دهز}} \textbf{\rzl{AbuAmr}{حطي}} \textbf{\rzl{IbnAmir}{كلم}} \textbf{\rzl{Asim}{نصع}} \textbf{\rzl{Hamza}{فضق}} \textbf{\rzl{Kisai}{رست}}\\
(الحرف الأول للإمام، والثاني لراويه الأول، والثالث لراويه الثاني.)
\end{tanbih}

\section{رموز الجماعة}
وهذه رموز لأكثر من قارئ، ولونها \textcolor{cGroup}{\textbf{بنيّ}}:

\begin{center}
\renewcommand{\arraystretch}{1.45}
\begin{tabular}{c l}
\toprule
\textbf{الرمز} & \textbf{القراء} \\
\midrule
\rzlg{ث} & الكوفيون: \Asim، \Hamza، \Kisai \\
\rzlg{خ} & السبعة إلا \Nafi \\
\rzlg{ذ} & \IbnAmir{} والكوفيون \\
\rzlg{ظ} & \IbnKathir{} والكوفيون \\
\rzlg{غ} & \AbuAmr{} والكوفيون \\
\rzlg{ش} & \Hamza{} و\Kisai \\
\midrule
\rzlg{صُحْبَة} & \Hamza، \Kisai، \Shuba \\
\rzlg{صِحَاب} & \Hamza، \Kisai، \Hafs \\
\rzlg{عَمّ} & \Nafi، \IbnAmir \\
\rzlg{سَمَا} & \Nafi، \IbnKathir، \AbuAmr \\
\rzlg{حَقّ} & \IbnKathir، \AbuAmr \\
\rzlg{نَفَر} & \IbnKathir، \AbuAmr، \IbnAmir \\
\rzlg{حِرْمِيّ} & \Nafi، \IbnKathir \\
\rzlg{حِصْن} & \Nafi{} والكوفيون \\
\bottomrule
\end{tabular}
\end{center}

\section{معاني الصناديق والألوان الأخرى}
\begin{center}
\renewcommand{\arraystretch}{1.5}
\begin{tabular}{c l}
\toprule
\textbf{الشكل} & \textbf{معناه} \\
\midrule
\colorbox{parchment}{\textcolor{gold}{\textbf{إطار ذهبي}}} & نص البيت من الشاطبية \\
\qw{كلمة مظللة} & الكلمة القرآنية المختلف فيها داخل البيت \\
\rz{Asim}{كلمة ملونة} & كلمة فيها رمز قارئ (لونها لون القارئ) \\
\colorbox{vocabback}{\textcolor{vocabframe}{\textbf{رمادي}}} & معاني الكلمات \\
\colorbox{white}{\textcolor{gold!80!black}{\textbf{ذهبي}}} & القراءات في البيت: من قرأ بماذا \\
\colorbox{ruleback}{\textcolor{ruleframe}{\textbf{أخضر}}} & قاعدة أو شرح اصطلاح \\
\colorbox{exback}{\textcolor{exframe}{\textbf{أزرق}}} & أمثلة من القرآن \\
\colorbox{tahrirback}{\textcolor{tahrirframe}{\textbf{بنفسجي}}} & تحريرات (من إتحاف البرية) \\
\colorbox{noteback}{\textcolor{noteframe}{\textbf{برتقالي}}} & تنبيه \\
\colorbox{sumback}{\textcolor{sumframe}{\textbf{خمري}}} & خلاصة الباب \\
\ayah{آية} & نص قرآني بالرسم العثماني، و\src{السورة: الآية} موضعه \\
\bottomrule
\end{tabular}
\end{center}

\begin{tanbih}[ملاحظة على التحريرات]
«التحريرات» هي ضبط ما يُقرأ به من طريق الشاطبية تحديدًا، ومنع خلط الطرق. وما نقلناه منها هو من منظومة «إتحاف البرية بتحريرات الشاطبية» للشيخ حسن خلف الحسيني رحمه الله. والأخذ بالتحريرات في الأداء يكون بالتلقي من شيخ متقن.
\end{tanbih}
